\section{Results and Analysis}
\label{ch:results}

% =============================================================================
% SKELETON --- restructured 2026-08-18 for the 10,000-word limit and the revised
% research questions; scenario recast 2026-09-10 (supervisor_revision_plan.md item 4).
% Written once the final run exists (20 replicates, shortfall borrowing path capped).
%
% BUDGET: 3,300 words of PROSE. Tables and figures do not count, so every number
% goes in a table and prose is spent only on what the number means. A results
% chapter that narrates its numbers runs to twice this length and reads worse.
%
% Section budgets, which are binding:
%   5.1 Baseline and RQ1   700
%   5.2 RQ2                1,100
%   5.3 RQ3                1,000
%   5.4 Robustness         500
%
% FIVE figures in the body, no more:
%   baseline_arrears_profile / rq1_stacking / rq2_default_surface /
%   rq3_scenarios_beta1 / robustness_tornado
% Everything else -> Appendix~\ref{app:supplementary}.
%
% Numbers below are from the 13 Aug working run (950 runs x 5 replicates,
% noise floor 0.33pp). They are placeholders: refresh every one against the
% final run before submission.
%
% EVERY table and figure carries a short \caption and a \tabnote that states: the
% replicate count and seed scheme; the noise floor (0.33pp; differences below
% ~0.67pp are noise); which scenario each column is; that access rates are shares
% of the BANKED subpopulation (82.8% ceiling), not of all households; and that
% beta = 0 is the control arm. This is a caption requirement, not a prose habit.
% =============================================================================

\subsection{Baseline and the Household Population}
\label{sec:results-baseline}

% BUDGET: 700 words. Answers the second half of RQ1.
%
% CONTENT REQUIRED:
% - The fitted income-shock probability and payment-inattention rate, and the arrears
%   profile they produce against Table~\ref{tab:ccmr}. State plainly that two parameters
%   were fitted to two bands, so those two bands prove nothing (Section~\ref{sec:lim-calibration}).
% - THE RESULT THAT MATTERS: the 60+ day band, 15.38% against 16.54%, never fitted.
% - Table: the seven independent checks (60+ arrears, 2+ facilities, BNPL spend per user,
%   average purchase, complementarity direction, DTI peak, zero savings). Six of seven
%   pass; zero savings is 45% against 36%, same order but high, and it rests on the
%   weakest variable in the data layer.
% - Foreground patterns 2 and 4 (D4). Patterns 1 and 3 get one line each in the same
%   table --- do not give them their own prose.
% - Close RQ1 explicitly: a population built this way does reproduce behaviour it was
%   not fitted to, on four of five independent behavioural checks, with the failure named.
% - Figure: baseline_arrears_profile.png

\subsection{RQ2: Stacking and Population Default}
\label{sec:results-rq2}

% BUDGET: 1,100 words. The largest section.
%
% CONTENT REQUIRED:
% - Stacking is EMERGENT, not programmed. 15-39% of adopters hold 2+ facilities against
%   the CFPB's 32%, as an order-of-magnitude check. Platform-count sweep 1-6; the
%   single-platform arm gives exactly zero stacking, which isolates cross-platform
%   accumulation from single-platform accumulation.
% - Default: 24.2% with BNPL off, 27.2% on, 29.1% on with strong peer influence.
%   The beta=0 row must be explicit in every table and on every figure.
% - Access sweep runs inside the banked subpopulation (82.8% ceiling, Submodel~12).
%   State the ceiling wherever an access rate is reported so it is not misread as a
%   share of all households.
%
% - THE NEGATIVE RESULT, reported as a titled finding, not buried (D3):
%   the increase is SMOOTH in access. No cliff edge anywhere on the surface. This
%   falsifies the pre-registered non-linearity hypothesis.
%   Then the pre-registered Granovetter threshold variant, run precisely because a
%   linear rule might give a linear answer for trivial reasons: adoption roughly
%   doubles (21% -> 42%) but default moves about half a percentage point, and the
%   adoption line is almost perfectly straight (R^2 = 0.9999). Two structurally
%   distinct social mechanisms, neither producing a jump.
%   The structural reason is the paragraph below (Section~\ref{sec:no-cascade}), which
%   is already written and must sit directly after the negative result.
%
% - THE SPLIT ANSWER, which must be reported honestly and in one place: the credit
%   limit binds on 54% of requests but barely moves default, because a constrained
%   household borrows less per purchase rather than defaulting less. Reported apart,
%   these two facts read as a contradiction.
% - Pattern 1 is tested here: enabling BNPL raises traditional servicing stress by
%   +2.9 to +3.5pp. Attach the caveat of Section~\ref{sec:lim-validation} --- the
%   direction is partly designed into the lender-choice rule, the magnitude is not.
% - Figures: rq1_stacking.png, rq2_default_surface.png
%   NOTE: figure FILENAMES keep the original RQ numbering and are deliberately frozen
%   (they are baked into results/raw/ labels and analysis.py load keys). rq1_stacking.png
%   belongs to THIS section --- thesis RQ2. See the numbering note at the top of
%   simulation/experiments.py.
% - Table: default rate by access x beta, with the beta=0 row first.
% - beta is RQ2's experimental axis: it is the strength of peer influence on BNPL
%   adoption, so sweeping it is what separates a social explanation from a purely
%   financial one. Under RQ3 it becomes Scenario 3, crossed with the other two.

% --- Written prose: the structural explanation of the negative result. Relocated from
% the former Discussion section (6.5) and to be kept directly after the result it
% explains when this section is drafted. ---
\label{sec:no-cascade}
The absence of a cliff edge is a property of a model built this way rather than a general finding.
Six mechanisms are excluded by design (Table~\ref{tab:design-concepts}), and the second of them, the absence of any channel by which one household's default affects another household's income, is the explanation of the negative result.
With no such channel, what the model produces is many correlated individual defaults rather than a cascade, and a threshold in the population default rate requires feedback in the outcome itself: households catch \ac{BNPL} from their neighbours, but they do not catch financial distress from them.
A successor design with an impairable lender or a default-to-income channel could plausibly produce a cliff edge, so the negative result is bounded rather than general.

\subsection{RQ3: Three Scenarios for the Lending Environment}
\label{sec:results-rq3}

% BUDGET: 1,000 words. The most novel result in the thesis.
%
% SCENARIO SET (supervisor decision 2026-09-09: three scenarios, not four levers).
% The run grid in experiments.py is already factorial (for b in (0.0, 1.0) crossed with
% every switch), so present it as it is, 2 x 3:
%
%                    | benchmark        | bureau visibility | mandatory screening
%   beta = 0         | pre-2026 status  | Scenario 1        | Scenario 2
%   beta > 0         | Scenario 3       | 1 x 3             | 2 x 3
%
% - Each cell gets the SAME outcome panel as the benchmark (default, arrears bands,
%   adoption, cumulative BNPL volume, by quintile), NOT a delta column against it.
% - First sentence: Scenario 1 IS the 2026 NCR reporting requirement, represented as
%   its affordability channel only; the scoring channel is out of scope (Submodel 10).
% - Scenario 3 asks whether the answers to 1 and 2 depend on whether adoption is
%   socially transmitted. Its beta = 0 row is the control of every comparison.
% - Working-run magnitudes, to be refreshed against the final run:
%     bureau visibility            ~1%, essentially none
%     mandatory affordability check  essentially none
% - The bureau-visibility contrast deserves its own paragraph: it is the comparison the
%   whole model was built to make, it holds every uncalibrated parameter fixed across
%   both arms, and it survived a complete rebuild of the two parameters underneath it.
%   The near-null result is a finding, not a failure --- and it is now an ex-ante
%   statement about a regulation with a commencement date, so state its scope
%   (affordability channel) in the same sentence as the number.
% - The defer-versus-desist test is CUMULATIVE BNPL VOLUME over the full horizon, not
%   the timing of purchases. Unchanged volume with a shifted time distribution means
%   agents defer; lower cumulative volume means they desist.
% - THE ARGUMENT must be grounded on RQ2 evidence that stays in the body, not on the
%   removed facility cap: the mechanism of harm is a household holding several
%   facilities at once, established by the emergent 2+ facility share against the
%   CFPB 32%, the single-platform arm producing exactly zero stacking, and default
%   rising with platform count. Both retained instruments act on the individual
%   agreement (telling the lender what is owed; testing one loan in isolation) and
%   neither acts on the number of them, so their null results CORROBORATE the
%   mechanism rather than establish it.
% - Cooling-off window and concurrent-facility cap are CUT from the body (decision 3,
%   2026-09-09). Their runs (rq3.parquet: rq3_kcool*, rq3_cap*) go to
%   Appendix~\ref{app:supplementary} as tables with no body prose.
% - Figure: rq3_interventions_beta1.png regenerated with the two cut levers removed
%   (the beta=0 panel goes to the appendix, with the beta=0 row retained in the body
%   table).

\subsection{Sensitivity and Robustness}
\label{sec:results-robustness}

% BUDGET: 500 words. One figure, one table, prose only for what moves.
%
% CONTENT REQUIRED:
% - Replicate count, seed scheme, noise floor (0.33pp on the working run; differences
%   below ~0.67pp are treated as noise). Total run count and wall-clock cost.
% - Figure: robustness_tornado.png. Each bar is how far default moves across the
%   plausible range of one parameter; red bars are rules with no citation.
% - The largest remaining sensitivity is the Submodel~4 borrowing-amount rule
%   (Section~\ref{sec:amount}). Report it plainly: two of the three specifications
%   agree to within the noise floor (28.91% and 28.24%) and the whole range comes from
%   the third and most generous. State whether the capped shortfall path (added after
%   the working run) narrows it.
% - Confirm empirically, rather than asserting, that the peer channel is insensitive to
%   activation order.
% - Everything else --- Sobol indices, activation-order re-runs, population-size
%   stability at 1,000/5,000/10,000, the minimum-payer sweep 0.20-0.40, and the full
%   platform-count detail --- goes to Appendix~\ref{app:supplementary} as tables, with
%   one sentence here saying so. Do not narrate them.
